\begin{table*}
    \caption{The prompt used in the full instantiation of ResearchAgent for problem identification.}
    \label{tab:prompt:problem}
    \vspace{-0.1in}
    \resizebox{0.95\textwidth}{!}{
    \renewcommand{\arraystretch}{1.1}
    \renewcommand{\tabcolsep}{2.5mm}
        \begin{tabular}{ll}
        \toprule
        \multicolumn{1}{p{.18\textwidth}}{\textbf{Types}} & \makecell{\multicolumn{1}{p{.92\textwidth}}{\textbf{Texts}}} \\
        \midrule
        \multicolumn{1}{p{.18\textwidth}}{\textbf{System Message}} & 
        \makecell{
            \multicolumn{1}{p{.92\textwidth}}{You are an AI assistant whose primary goal is to identify promising, new, and key scientific problems based on existing scientific literature, in order to aid researchers in discovering novel and significant research opportunities that can advance the field.}
        } \\
        \midrule
        \multicolumn{1}{p{.18\textwidth}}{\textbf{User Message}} & 
        \makecell{
            \multicolumn{1}{p{.92\textwidth}}{{You are going to generate a research problem that should be original, clear, feasible, relevant, and significant to its field. This will be based on the title and abstract of the target paper, those of \{len(references)\} related papers in the existing literature, and \{len(entities)\} entities potentially connected to the research area.}} \\ \\
            \multicolumn{1}{p{.92\textwidth}}{{Understanding of the target paper, related papers, and entities is essential:}} \\
            \multicolumn{1}{p{.92\textwidth}}{{- The target paper is the primary research study you aim to enhance or build upon through future research, serving as the central source and focus for identifying and developing the specific research problem.}} \\
            \multicolumn{1}{p{.92\textwidth}}{{- The related papers are studies that have cited the target paper, indicating their direct relevance and connection to the primary research topic you are focusing on, and providing additional context and insights that are essential for understanding and expanding upon the target paper.}} \\
            \multicolumn{1}{p{.92\textwidth}}{{- The entities can include topics, keywords, individuals, events, or any subjects with possible direct or indirect connections to the target paper or the related studies, serving as auxiliary sources of inspiration or information that may be instrumental in formulating the research problem.}} \\ \\
            \multicolumn{1}{p{.92\textwidth}}{{Your approach should be systematic:}} \\ 
            \multicolumn{1}{p{.92\textwidth}}{{- Start by thoroughly reading the title and abstract of the target paper to understand its core focus.}} \\
            \multicolumn{1}{p{.92\textwidth}}{{- Next, proceed to read the titles and abstracts of the related papers to gain a broader perspective and insights relevant to the primary research topic.}} \\
            \multicolumn{1}{p{.92\textwidth}}{{- Finally, explore the entities to further broaden your perspective, drawing upon a diverse pool of inspiration and information, while keeping in mind that not all may be relevant.}} \\ \\
            \multicolumn{1}{p{.92\textwidth}}{{I am going to provide the target paper, related papers, and entities, as follows: }} \\
            \multicolumn{1}{p{.92\textwidth}}{{Target paper title: \{paper['title']\}}} \\
            \multicolumn{1}{p{.92\textwidth}}{{Target paper abstract: \{paper['abstract']\} }} \\
            \multicolumn{1}{p{.92\textwidth}}{{Related paper titles: \{relatedPaper['titles']\}}} \\
            \multicolumn{1}{p{.92\textwidth}}{{Related paper abstracts: \{relatedPaper['abstracts']\} }} \\
            \multicolumn{1}{p{.92\textwidth}}{{Entities: \{Entities\} }} \\ \\
            \multicolumn{1}{p{.92\textwidth}}{{With the provided target paper, related papers, and entities, your objective now is to formulate a research problem that not only builds upon these existing studies but also strives to be original, clear, feasible, relevant, and significant. Before crafting the research problem, revisit the title and abstract of the target paper, to ensure it remains the focal point of your research problem identification process.}} \\ \\
            \multicolumn{1}{p{.92\textwidth}}{{Target paper title: \{paper['title']\}}} \\
            \multicolumn{1}{p{.92\textwidth}}{{Target paper abstract: \{paper['abstract']\} }} \\ \\
            \multicolumn{1}{p{.92\textwidth}}{{Then, following your review of the above content, please proceed to generate one research problem with the rationale, in the format of }} \\
            \multicolumn{1}{p{.92\textwidth}}{{Problem: }} \\
            \multicolumn{1}{p{.92\textwidth}}{{Rationale: }} \\
        } \\
        \bottomrule
        \end{tabular}
    }
\end{table*}

\begin{table*}
    \caption{The prompt used in the full instantiation of ResearchAgent for method development.}
    \label{tab:prompt:method}
    \vspace{-0.1in}
    \resizebox{0.95\textwidth}{!}{
    \renewcommand{\arraystretch}{1.1}
    \renewcommand{\tabcolsep}{2.5mm}
        \begin{tabular}{ll}
        \toprule
        \multicolumn{1}{p{.18\textwidth}}{\textbf{Types}} & \makecell{\multicolumn{1}{p{.92\textwidth}}{\textbf{Texts}}} \\
        \midrule
        \multicolumn{1}{p{.18\textwidth}}{\textbf{System Message}} & 
        \makecell{
            \multicolumn{1}{p{.92\textwidth}}{You are an AI assistant whose primary goal is to propose innovative, rigorous, and valid methodologies to solve newly identified scientific problems derived from existing scientific literature, in order to empower researchers to pioneer groundbreaking solutions that catalyze breakthroughs in their fields.}
        } \\
        \midrule
        \multicolumn{1}{p{.18\textwidth}}{\textbf{User Message}} & 
        \makecell{
            \multicolumn{1}{p{.92\textwidth}}{{You are going to propose a scientific method to address a specific research problem. Your method should be clear, innovative, rigorous, valid, and generalizable. This will be based on a deep understanding of the research problem, its rationale, existing studies, and various entities.}} \\ \\
            \multicolumn{1}{p{.92\textwidth}}{{Understanding of the research problem, existing studies, and entities is essential:}} \\
            \multicolumn{1}{p{.92\textwidth}}{{- The research problem has been formulated based on an in-depth review of existing studies and a potential exploration of relevant entities, which should be the cornerstone of your method development.}} \\
            \multicolumn{1}{p{.92\textwidth}}{{- The existing studies refer to the target paper that has been pivotal in identifying the problem, as well as the related papers that have been additionally referenced in the problem discovery phase, all serving as foundational material for developing the method.}} \\
            \multicolumn{1}{p{.92\textwidth}}{{- The entities can include topics, keywords, individuals, events, or any subjects with possible direct or indirect connections to the existing studies, serving as auxiliary sources of inspiration or information that may be instrumental in method development.}} \\ \\
            \multicolumn{1}{p{.92\textwidth}}{{Your approach should be systematic:}} \\ 
            \multicolumn{1}{p{.92\textwidth}}{{- Start by thoroughly reading the research problem and its rationale, to understand your primary focus.}} \\
            \multicolumn{1}{p{.92\textwidth}}{{- Next, proceed to review the titles and abstracts of existing studies, to gain a broader perspective and insights relevant to the primary research topic.}} \\
            \multicolumn{1}{p{.92\textwidth}}{{- Finally, explore the entities to further broaden your perspective, drawing upon a diverse pool of inspiration and information, while keeping in mind that not all may be relevant.}} \\ \\
            \multicolumn{1}{p{.92\textwidth}}{{I am going to provide the research problem, existing studies (target paper \& related papers), and entities, as follows: }} \\
            \multicolumn{1}{p{.92\textwidth}}{{Research problem: \{researchProblem\}}} \\
            \multicolumn{1}{p{.92\textwidth}}{{Rationale: \{researchProblemRationale\}}} \\
            \multicolumn{1}{p{.92\textwidth}}{{Target paper title: \{paper['title']\}}} \\
            \multicolumn{1}{p{.92\textwidth}}{{Target paper abstract: \{paper['abstract']\} }} \\
            \multicolumn{1}{p{.92\textwidth}}{{Related paper titles: \{relatedPaper['titles']\}}} \\
            \multicolumn{1}{p{.92\textwidth}}{{Related paper abstracts: \{relatedPaper['abstracts']\} }} \\
            \multicolumn{1}{p{.92\textwidth}}{{Entities: \{Entities\} }} \\ \\
            \multicolumn{1}{p{.92\textwidth}}{{With the provided research problem, existing studies, and entities, your objective now is to formulate a method that not only leverages these resources but also strives to be clear, innovative, rigorous, valid, and generalizable. Before crafting the method, revisit the research problem, to ensure it remains the focal point of your method development process.}} \\ \\
            \multicolumn{1}{p{.92\textwidth}}{{Research problem: \{researchProblem\}}} \\
            \multicolumn{1}{p{.92\textwidth}}{{Rationale: \{researchProblemRationale\}}} \\ \\
            \multicolumn{1}{p{.92\textwidth}}{{Then, following your review of the above content, please proceed to propose your method with its rationale, in the format of }} \\
            \multicolumn{1}{p{.92\textwidth}}{{Method: }} \\
            \multicolumn{1}{p{.92\textwidth}}{{Rationale: }} \\
        } \\
        \bottomrule
        \end{tabular}
    }
\end{table*}

\begin{table*}
    \caption{The prompt used in the full instantiation of ResearchAgent for experiment design.}
    \label{tab:prompt:experiment}
    \vspace{-0.1in}
    \resizebox{0.95\textwidth}{!}{
    \renewcommand{\arraystretch}{1.1}
    \renewcommand{\tabcolsep}{2.5mm}
        \begin{tabular}{ll}
        \toprule
        \multicolumn{1}{p{.18\textwidth}}{\textbf{Types}} & \makecell{\multicolumn{1}{p{.92\textwidth}}{\textbf{Texts}}} \\
        \midrule
        \multicolumn{1}{p{.18\textwidth}}{\textbf{System Message}} & 
        \makecell{
            \multicolumn{1}{p{.92\textwidth}}{You are an AI assistant whose primary goal is to design robust, feasible, and impactful experiments based on identified scientific problems and proposed methodologies from existing scientific literature, in order to enable researchers to systematically test hypotheses and validate groundbreaking discoveries that can transform their respective fields.}
        } \\
        \midrule
        \multicolumn{1}{p{.18\textwidth}}{\textbf{User Message}} & 
        \makecell{
            \multicolumn{1}{p{.92\textwidth}}{{You are going to design an experiment, aimed at validating a proposed method to address a specific research problem. Your experiment design should be clear, robust, reproducible, valid, and feasible. This will be based on a deep understanding of the research problem, scientific method, existing studies, and various entities.}} \\ \\
            \multicolumn{1}{p{.92\textwidth}}{{Understanding of the research problem, scientific method, existing studies, and entities is essential:}} \\
            \multicolumn{1}{p{.92\textwidth}}{{- The research problem has been formulated based on an in-depth review of existing studies and a potential exploration of relevant entities.}} \\
            \multicolumn{1}{p{.92\textwidth}}{{- The scientific method has been proposed to tackle the research problem, which has been informed by insights gained from existing studies and relevant entities.}} \\
            \multicolumn{1}{p{.92\textwidth}}{{- The existing studies refer to the target paper that has been pivotal in identifying the problem and method, as well as the related papers that have been additionally referenced in the discovery phase of the problem and method, all serving as foundational material for designing the experiment.}} \\
            \multicolumn{1}{p{.92\textwidth}}{{- The entities can include topics, keywords, individuals, events, or any subjects with possible direct or indirect connections to the existing studies, serving as auxiliary sources of inspiration or information that may be instrumental in your experiment design.}} \\ \\
            \multicolumn{1}{p{.92\textwidth}}{{Your approach should be systematic:}} \\ 
            \multicolumn{1}{p{.92\textwidth}}{{- Start by thoroughly reading the research problem and its rationale followed by the proposed method and its rationale, to pinpoint your primary focus.}} \\
            \multicolumn{1}{p{.92\textwidth}}{{- Next, proceed to review the titles and abstracts of existing studies, to gain a broader perspective and insights relevant to the primary research topic.}} \\
            \multicolumn{1}{p{.92\textwidth}}{{- Finally, explore the entities to further broaden your perspective, drawing upon a diverse pool of inspiration and information, while keeping in mind that not all may be relevant.}} \\ \\
            \multicolumn{1}{p{.92\textwidth}}{{I am going to provide the research problem, scientific method, existing studies (target paper \& related papers), and entities, as follows: }} \\
            \multicolumn{1}{p{.92\textwidth}}{{Research problem: \{researchProblem\}}} \\
            \multicolumn{1}{p{.92\textwidth}}{{Rationale: \{researchProblemRationale\}}} \\
            \multicolumn{1}{p{.92\textwidth}}{{Scientific method: \{scientificMethod\}}} \\
            \multicolumn{1}{p{.92\textwidth}}{{Rationale: \{scientificMethodRationale\}}} \\
            \multicolumn{1}{p{.92\textwidth}}{{Target paper title: \{paper['title']\}}} \\
            \multicolumn{1}{p{.92\textwidth}}{{Target paper abstract: \{paper['abstract']\} }} \\
            \multicolumn{1}{p{.92\textwidth}}{{Related paper titles: \{relatedPaper['titles']\}}} \\
            \multicolumn{1}{p{.92\textwidth}}{{Related paper abstracts: \{relatedPaper['abstracts']\} }} \\
            \multicolumn{1}{p{.92\textwidth}}{{Entities: \{Entities\} }} \\ \\
            \multicolumn{1}{p{.92\textwidth}}{{With the provided research problem, scientific method, existing studies, and entities, your objective now is to design an experiment that not only leverages these resources but also strives to be clear, robust, reproducible, valid, and feasible. Before crafting the experiment design, revisit the research problem and proposed method, to ensure they remain at the center of your experiment design process.}} \\ \\
            \multicolumn{1}{p{.92\textwidth}}{{Research problem: \{researchProblem\}}} \\
            \multicolumn{1}{p{.92\textwidth}}{{Rationale: \{researchProblemRationale\}}} \\ 
            \multicolumn{1}{p{.92\textwidth}}{{Scientific method: \{scientificMethod\}}} \\
            \multicolumn{1}{p{.92\textwidth}}{{Rationale: \{scientificMethodRationale\}}} \\ \\
            \multicolumn{1}{p{.92\textwidth}}{{Then, following your review of the above content, please proceed to outline your experiment with its rationale, in the format of }} \\
            \multicolumn{1}{p{.92\textwidth}}{{Experiment: }} \\
            \multicolumn{1}{p{.92\textwidth}}{{Rationale: }} \\
        } \\
        \bottomrule
        \end{tabular}
    }
\end{table*}